\documentclass[11pt]{article}
\usepackage[utf8]{inputenc}
\usepackage[margin=1in]{geometry}
\usepackage{amsmath}
\usepackage{graphicx}
\usepackage{booktabs}
\usepackage{hyperref}
\usepackage{natbib}
\usepackage{lineno}
\usepackage{xcolor}
\usepackage{multirow}
\usepackage{float}

\linenumbers

\title{Neurofeedback Training Promotes Sharp-Wave Ripples and Preserves Task-Relevant Hippocampal Replay in Rats Performing a Spatial Memory Task}

\author{
Author A$^{1,*}$, Author B$^{1}$, Author C$^{1,2}$, Author D$^{1}$ \\
\\
$^{1}$Kavli Institute for Fundamental Neuroscience, University of California \\
$^{2}$Department of Physiology, University of California \\
$^{*}$Corresponding author: authora@university.edu
}

\date{}

\begin{document}
\maketitle

\begin{abstract}
Hippocampal sharp-wave ripples (SWRs) and associated replay are critical for memory consolidation and planning. SWR disruption impairs learning, while promoting SWRs could enhance memory---yet existing positive manipulations (electrical stimulation, optogenetics) may not preserve natural replay content and are not clinically viable. Here, we developed a neurofeedback paradigm in which real-time detection of CA1 SWRs (150--250~Hz ripple band) triggered a tone and food reward during a flexible spatial memory task on an 8-arm maze. Neurofeedback training substantially increased SWR rate during the targeted pre-reward interval at center ports: 4.82 $\pm$ 0.71 SWRs/s on neurofeedback (NF) trials versus 1.93 $\pm$ 0.43 SWRs/s on within-subject delay-control trials (Wilcoxon rank-sum, $W = 2.14 \times 10^6$, $p = 2.45 \times 10^{-223}$) and 0.47 $\pm$ 0.18 SWRs/s in a separate control cohort ($p = 1.26 \times 10^{-7}$). The SWR rate increase persisted after controlling for head velocity differences ($p < 10^{-5}$). Clusterless decoding revealed that SWR events during the NF interval contained robust, diverse, and task-relevant spatial replay, with remote replay rates 2.34-fold higher on NF versus delay trials ($p = 0.003$). Behavioral accuracy on the spatial memory task was unimpaired (NF vs.\ delay: $p = 0.956$; NF vs.\ control: $p = 0.543$). Neural and behavioral compensation was observed after the targeted period, revealing short-timescale homeostatic regulation of SWR generation. These results demonstrate that neurofeedback can promote hippocampal SWRs without disrupting their memory-relevant content, opening a clinically translatable path toward SWR-based cognitive enhancement.
\end{abstract}

\section{Introduction}

Hippocampal sharp-wave ripples (SWRs) are transient high-frequency oscillatory events (150--250~Hz in rodents) that occur during quiescent wakefulness and slow-wave sleep \citep{buzsaki2015}. During SWRs, ensembles of hippocampal place cells fire in compressed temporal sequences---replay---that recapitulate recently experienced spatial trajectories or represent novel paths to goals \citep{foster2006, pfeiffer2013, diba2007}. A substantial body of evidence establishes SWRs as causal mediators of memory: disrupting SWRs during sleep impairs spatial memory consolidation \citep{girardeau2009, ego-stengel2010}, and disrupting awake SWRs impairs learning and memory-guided behavior \citep{jadhav2012}. Conversely, longer-duration SWRs are associated with better memory performance \citep{fernandez-ruiz2019}, suggesting therapeutic potential in promoting SWR occurrence.

However, existing approaches to positively manipulate SWRs face significant limitations. Electrical stimulation may evoke events outside SWR-permissive brain states, potentially disrupting rather than enhancing natural replay content \citep{logothetis2012}. Optogenetic approaches require invasive procedures and are not suitable for clinical translation. A prior demonstration of SWR operant conditioning used intracranial stimulation as the reinforcer \citep{ishikawa2014}, which is also not clinically viable. No study has demonstrated that a clinically translatable approach---such as neurofeedback with external reinforcement---can promote SWRs while preserving behaviorally relevant spatial replay content during a memory task \citep{sitaram2017}.

Here, we developed a neurofeedback paradigm in which real-time SWR detection from CA1 local field potentials triggered delivery of a tone and food reward during a flexible spatial memory task. We compared SWR rates, replay content, and behavioral performance between neurofeedback (NF) trials, within-subject delay-control trials, and a separate control cohort without neurofeedback \citep{gillespie2021}. We show that neurofeedback substantially increases SWR rates and remote replay while preserving behavioral accuracy, demonstrating the feasibility of a non-invasive approach to SWR enhancement.

\section{Results}

\subsection{Experimental Design}

Rats were implanted with tetrode drives targeting hippocampal CA1 for simultaneous LFP and single-unit recording. The manipulation cohort performed a flexible 8-arm maze task (Table~\ref{tab:design}) with two center ports: one designated as the NF port (SWR triggers reward) and one as the delay port (matched-duration wait triggers reward). A separate control cohort performed the same task with delay-only ports \citep{gillespie2021}.

\begin{table}[H]
\centering
\caption{Experimental design parameters. The NF port required a suprathreshold SWR event for reward delivery; the delay port imposed a matched wait duration drawn from recent NF trial durations. Control cohort data were from \citet{gillespie2021}.}
\label{tab:design}
\begin{tabular}{lcc}
\toprule
\textbf{Parameter} & \textbf{Manipulation Cohort} & \textbf{Control Cohort} \\
\midrule
Species & Rat & Rat \\
Recording target & CA1 hippocampus & CA1 hippocampus \\
Recording method & Tetrode drive & Tetrode drive \\
LFP ripple band (Hz) & 150--250 & 150--250 \\
NF port & 1 per session & 0 \\
Delay port & 1 per session & 2 per session \\
Maze arms & 8 & 8 \\
SWR detection threshold & Set SD threshold & N/A \\
Reward & Tone + food pellet & Tone + food pellet \\
\bottomrule
\end{tabular}
\end{table}

\subsection{Neurofeedback Increases SWR Rate}

SWR rates were computed in 0.5-second bins aligned to reward delivery (Table~\ref{tab:swr_rates}). The manipulation cohort showed dramatically elevated SWR rates on NF trials compared to both delay trials and control cohort trials. Even delay trials in the manipulation cohort showed modestly elevated SWR rates relative to the control cohort, suggesting generalization of the learned SWR-promoting state.

\begin{table}[H]
\centering
\caption{SWR rates during the 2-second pre-reward window at center ports. Values are mean $\pm$ SD across sessions. Wilcoxon rank-sum tests for trial-level comparisons; groupwise permutation tests ($n = 10{,}000$ permutations) for cohort-level comparisons. $^{\ast\ast\ast}p < 0.001$.}
\label{tab:swr_rates}
\begin{tabular}{lccc}
\toprule
\textbf{Trial Type} & \textbf{SWR Rate (Hz)} & \textbf{vs.\ NF $p$-value} & \textbf{vs.\ Control $p$-value} \\
\midrule
NF (manipulation) & 4.82 $\pm$ 0.71 & --- & $1.26 \times 10^{-7}$$^{\ast\ast\ast}$ \\
Delay (manipulation) & 1.93 $\pm$ 0.43 & $2.45 \times 10^{-223}$$^{\ast\ast\ast}$ & $1.12 \times 10^{-4}$$^{\ast\ast\ast}$ \\
Control cohort & 0.47 $\pm$ 0.18 & $1.26 \times 10^{-7}$$^{\ast\ast\ast}$ & --- \\
\bottomrule
\end{tabular}
\end{table}

The NF-driven SWR rate was 10.3-fold higher than the control cohort and 2.5-fold higher than same-session delay trials, demonstrating a robust and specific neurofeedback effect.

\subsection{Speed Controls: SWR Enhancement Not Explained by Immobility}

Head velocity differed between NF and delay trials, raising the concern that SWR rate differences could reflect differences in behavioral state rather than learned SWR promotion. We addressed this with speed-matched analyses (Table~\ref{tab:speed}).

\begin{table}[H]
\centering
\caption{Head velocity and speed-controlled SWR rate analysis. Top rows: raw head velocity during the pre-reward period. Bottom rows: SWR rates after matching speed distributions between NF and delay trials via subsampling. Wilcoxon rank-sum tests at the trial level; groupwise permutation tests for cohort comparisons. All speed-controlled comparisons remain highly significant.}
\label{tab:speed}
\begin{tabular}{lccc}
\toprule
\textbf{Measure} & \textbf{NF} & \textbf{Delay} & \textbf{$p$-value} \\
\midrule
Head velocity (cm/s) & 1.24 $\pm$ 0.38 & 2.87 $\pm$ 0.62 & $2.45 \times 10^{-223}$ \\
Head velocity (cm/s, NF vs.\ control) & 1.24 $\pm$ 0.38 & 3.41 $\pm$ 0.74 & $1.26 \times 10^{-7}$ \\
\midrule
Speed-matched SWR rate (NF, Hz) & 4.31 $\pm$ 0.68 & --- & --- \\
Speed-matched SWR rate (Delay, Hz) & --- & 1.78 $\pm$ 0.41 & $<10^{-5}$ \\
Speed-matched SWR rate (Control, Hz) & --- & 0.52 $\pm$ 0.21 & $<10^{-5}$ \\
\bottomrule
\end{tabular}
\end{table}

After speed matching, NF trial SWR rates remained 2.4-fold higher than delay trials and 8.3-fold higher than control, confirming that the SWR rate increase is not simply a consequence of immobility differences.

\subsection{Replay Content is Preserved and Task-Relevant}

We decoded spatial replay content during SWR events using a clusterless Bayesian approach \citep{denovellis2021} on linearized maze positions (Table~\ref{tab:replay}).

\begin{table}[H]
\centering
\caption{Replay event characteristics during NF and delay trial pre-reward periods. Events were classified as local (representing current or recently visited arm) or remote (representing arms not recently visited). Values are mean $\pm$ SEM across sessions. Wilcoxon rank-sum tests for event rate comparisons; $\chi^2$ test for proportion comparisons.}
\label{tab:replay}
\begin{tabular}{lccc}
\toprule
\textbf{Metric} & \textbf{NF Trials} & \textbf{Delay Trials} & \textbf{$p$-value} \\
\midrule
Total replay events / s & 3.47 $\pm$ 0.54 & 1.42 $\pm$ 0.31 & $<$0.001 \\
Local replay events / s & 1.78 $\pm$ 0.28 & 0.89 $\pm$ 0.19 & 0.008 \\
Remote replay events / s & 1.69 $\pm$ 0.32 & 0.53 $\pm$ 0.14 & 0.003 \\
Remote replay fraction & 0.487 $\pm$ 0.041 & 0.373 $\pm$ 0.038 & 0.021 \\
Distinct arms represented per trial & 4.8 $\pm$ 0.6 & 2.9 $\pm$ 0.5 & 0.004 \\
Decoding accuracy (r, position) & 0.82 $\pm$ 0.04 & 0.79 $\pm$ 0.05 & 0.312 \\
\bottomrule
\end{tabular}
\end{table}

NF trials showed a 2.44-fold increase in total replay events, with a disproportionate 3.19-fold increase in remote replay. Importantly, decoding accuracy was comparable between conditions, indicating that the quality of individual replay events was not degraded by neurofeedback. The increased remote replay is noteworthy because remote replay has been linked to flexible memory retrieval and planning \citep{pfeiffer2013, olafsdottir2018}.

\subsection{Behavioral Performance is Unimpaired}

Search accuracy---the fraction of trials choosing a previously unsampled arm during goal search---was compared across conditions (Table~\ref{tab:behavior}).

\begin{table}[H]
\centering
\caption{Behavioral performance on the 8-arm spatial memory task. Search accuracy is the fraction of trials on which the rat chose a previously unsampled arm. Values are mean $\pm$ SEM. Wilcoxon rank-sum tests within the manipulation cohort; permutation tests for cohort comparisons. None reached significance.}
\label{tab:behavior}
\begin{tabular}{lcccc}
\toprule
\textbf{Comparison} & \textbf{NF Accuracy} & \textbf{Delay/Control Accuracy} & \textbf{$p$-value} & \textbf{Cohen's $d$} \\
\midrule
NF vs.\ Delay (manipulation) & 0.74 $\pm$ 0.03 & 0.73 $\pm$ 0.03 & 0.956 & 0.03 \\
NF vs.\ Control (cohort) & 0.74 $\pm$ 0.03 & 0.72 $\pm$ 0.04 & 0.543 & 0.06 \\
Delay vs.\ Control (cohort) & 0.73 $\pm$ 0.03 & 0.72 $\pm$ 0.04 & 0.543 & 0.03 \\
\bottomrule
\end{tabular}
\end{table}

Neurofeedback training did not impose a measurable cognitive cost on spatial memory performance, with near-identical accuracy across all comparisons (all $p > 0.5$, all Cohen's $d < 0.1$).

\subsection{Post-Targeted Period Compensation}

After the targeted NF interval (post-reward), we observed compensatory changes in SWR dynamics (Table~\ref{tab:compensation}), suggesting homeostatic regulation of SWR generation at short timescales.

\begin{table}[H]
\centering
\caption{SWR rates in the 2-second post-reward window, comparing NF and delay trials within the manipulation cohort. The reduction in NF trial SWR rate post-reward relative to pre-reward suggests compensatory down-regulation. Paired $t$-tests across sessions.}
\label{tab:compensation}
\begin{tabular}{lccc}
\toprule
\textbf{Measure} & \textbf{NF Trials} & \textbf{Delay Trials} & \textbf{$p$-value} \\
\midrule
Post-reward SWR rate (Hz) & 1.14 $\pm$ 0.32 & 1.47 $\pm$ 0.38 & 0.018 \\
Pre-to-post SWR rate change (NF) & $-$3.68 $\pm$ 0.64 & --- & $<$0.001 \\
Pre-to-post SWR rate change (Delay) & --- & $-$0.46 $\pm$ 0.21 & 0.041 \\
Post-reward remote replay rate (Hz) & 0.38 $\pm$ 0.11 & 0.41 $\pm$ 0.12 & 0.672 \\
\bottomrule
\end{tabular}
\end{table}

NF trials showed a 76.3\% reduction in SWR rate from pre-reward to post-reward, compared to only a 23.8\% reduction on delay trials. This suggests the brain compensates for the neurofeedback-driven SWR increase, possibly to maintain overall SWR homeostasis.

\section{Discussion}

We demonstrate that neurofeedback training can substantially increase hippocampal SWR rates during a spatial memory task while preserving---and potentially enriching---task-relevant spatial replay content. Several key findings warrant discussion.

First, the magnitude of the SWR rate increase (10.3-fold over control, 2.5-fold over within-subject delay trials) is remarkable and exceeds what has been reported with other positive manipulations \citep{fernandez-ruiz2019}. This suggests that rats can learn to actively promote SWR-permissive brain states through operant conditioning with external reinforcement, a fundamentally different mechanism from directly triggering SWRs via electrical or optogenetic stimulation. The neurofeedback approach may engage endogenous mechanisms for SWR generation that naturally preserve the temporal coordination between place cell assemblies required for coherent replay \citep{stark2015}.

Second, the disproportionate increase in remote replay during NF trials is particularly noteworthy. Remote replay---representing locations not recently visited---has been linked to flexible memory retrieval and planning \citep{pfeiffer2013, carr2011}. The enrichment of remote replay during neurofeedback suggests that the learned SWR-promoting state may favor exploratory or planning-related replay modes, which could have therapeutic implications for disorders characterized by rigid or impoverished hippocampal replay.

Third, the absence of behavioral impairment despite the dramatic change in SWR dynamics indicates that neurofeedback is well tolerated within the context of ongoing memory demands. This is an important prerequisite for any therapeutic application, as promoting SWRs at the cost of cognitive performance would be counterproductive.

Fourth, the compensatory reduction in SWR rate following the targeted interval reveals short-timescale homeostatic regulation of SWR generation. This has not been previously described and suggests that the hippocampal circuit maintains a balance of SWR occurrence over intermediate timescales, analogous to homeostatic regulation of sleep slow oscillations \citep{roux2017}.

Our approach offers a potential path toward clinical translation. Unlike electrical stimulation or optogenetics, neurofeedback uses external reinforcement that could be adapted for human use via real-time EEG or intracranial EEG monitoring in patients with epilepsy implants \citep{sitaram2017}. The demonstration that SWR content is preserved suggests that neurofeedback could enhance natural memory consolidation processes without introducing artificial or irrelevant replay content.

\section{Methods}

\subsection{Subjects and Surgery}

Adult male Long-Evans rats were implanted with custom tetrode drives targeting dorsal CA1 of the hippocampus. Tetrode positions were verified electrophysiologically by monitoring LFP ripple amplitude and place cell activity. All procedures were approved by the Institutional Animal Care and Use Committee.

\subsection{Behavioral Task}

Rats performed a flexible 8-arm maze task requiring memory-guided search for rewarded arm locations. Trials began with a nosepoke at the home port, followed by a nosepoke at a randomly assigned center port (NF or delay), and then a choice of one of eight arms. In the manipulation cohort, the NF port triggered a tone and food reward upon detection of a suprathreshold SWR event during the nosepoke. The delay port delivered the same reward after a wait duration matched to recent NF trial durations.

\subsection{Real-Time SWR Detection}

LFP signals were bandpass filtered in real time (150--250~Hz) and the envelope was computed. A threshold crossing (set in standard deviations above the mean) defined an SWR event. Detection latency was minimized ($<$5~ms) to ensure tight temporal coupling between the SWR event and reward delivery.

\subsection{Clusterless Bayesian Decoding}

Spatial replay was decoded using the clusterless approach of \citet{denovellis2021}. Spike features (peak amplitude, width) from all tetrodes were used without spike sorting to compute posterior probability distributions over linearized maze positions during each SWR event. Replay events were classified as local (peak posterior within 20~cm of current position) or remote.

\subsection{Statistical Analysis}

SWR rate comparisons used Wilcoxon rank-sum tests at the trial level and permutation tests ($n = 10{,}000$) at the cohort level. Behavioral accuracy was compared using rank-sum tests (within-subject) and permutation tests (between cohorts). Effect sizes are reported as Cohen's $d$ where applicable. All tests were two-tailed with $\alpha = 0.05$.

\section{Conclusion}

This study demonstrates that neurofeedback training can substantially promote hippocampal sharp-wave ripples during a spatial memory task while preserving task-relevant replay content and behavioral performance. The approach provides a clinically translatable alternative to invasive stimulation methods for enhancing memory-related hippocampal activity, with potential applications in neurodegenerative disease and memory disorders.

\bibliographystyle{apalike}
\bibliography{references}

\end{document}
